上一单元确认了：条件合适时，初值问题有唯一解，解是一族互不相交的曲线。这一单元换一个问题——不写出解的公式，能否说出解的长期去向？答案是把解当作状态空间中的几何对象来研究：轨迹流向哪里、在哪里停住、绕什么打转，这些问题可以不依赖公式回答。

\begin{enumerate}
\def\labelenumi{\arabic{enumi}.}
\tightlist
\item
  相平面：把解画成几何：二维自治系统的向量场与轨迹，不动点按线性化特征值分类，极限环作为系统自产的周期行为。
\item
  稳定性与 Lyapunov 函数：线性化失效的临界情形下，用一个沿轨迹单调下降的``能量''函数证明稳定，并把梯度流与优化算法接上。
\item
  离散动力系统、分岔与混沌：迭代映射的不动点判据，Logistic 映射从稳定到倍周期再到混沌的完整剧情，以及对初值敏感为何让长期预测失效。
\end{enumerate}

三节依次放弃的东西越来越多：第一节还要向量场的解析表达式，第二节只需一个能量函数，第三节只剩一条迭代规则——但保留的始终是同一个判断：系统长期去向哪里。
